\section{Penalized reach and isolated interfaces}\label{sec:reach}

We now work in the only case left by finite-density rigidity: an
entire stable solution $v:\R^7\to(-1,1)$ with $D>0$ and unbounded
energy density.  Thus the interface dimension is $m=6$.
Our aim is to attach a length to each regular zero that controls both
bending of its own sheet and approach to another sheet.  We will prove
a differential inequality for this length, including a precise rule
for integrating the contribution of a second contact point.

Recall that $\Sigma=\{v=0\}$ and, on its regular part,
$\nu=\nabla U/|\nabla U|$, $A=d\nu|_{T\Sigma}$, $a=|A|$, and
$H=\tr A$.  For the opposite orientation, or a separately specified
unit normal $N$, we write $A_N$ and $H_N$.
The normal map is
\[
 (z,t)\longmapsto z+t\nu(z).
\]
It can lose local invertibility when $\Id+tA$ becomes singular;
even before that, normals from different parts of $\Sigma$ can meet.
The reach below records both obstructions through chords between zeros.

There is still useful local regularity in this setting.
Corollary~\ref{cor:finite-subcritical} applies whenever one large ball
has a fixed finite density bound; it does not require a density bound
on the whole solution.  We first find such a ball at almost the reach
scale.  Its isolated graph supports two parity projections that improve
mean-curvature derivatives.  A separate tangent-ball argument controls
the signed mean curvature at a possibly distant contact.  These are
the two local inputs to the reach inequality.

Throughout, put
\[
 K(R)=1+\sup_{x\in\R^7}M_R(x).
\]
Monotonicity and the uniform energy-integrand bound give
$1\le K(R)\le1+CR$ and $K(R)\to\infty$.
All constants are uniform in the solution and in translating centers.
When dependence on a fixed penalty matters, it is indicated by
$C_\Gamma$.  Dependence on $K$ will always be displayed explicitly.

\subsection{Chords, the penalty, and ordinary tubes}

For distinct zeros $z,y$, with $z$ regular, set
$\ell=|y-z|$, $e=(y-z)/\ell$, and $s=|\nu(z)\cdot e|$.
If $s>0$, then $\ell/s$ is twice the radius of the normal tangent
ball at $z$ passing through $y$.  As $y$ tends to $z$ along a
principal direction, the same quotient tends to twice the reciprocal
absolute principal curvature when that curvature is nonzero.
Thus taking an infimum over chords
measures both approach and bending.


Figure~\ref{fig:reach-tangent-ball} illustrates this quotient before
we modify its angular factor.

\begin{figure}[ht]
\centering
\begin{tikzpicture}[x=1.25cm,y=1.25cm,>=Latex,font=\small]
\draw[gray!65] (0,1.5) circle[radius=1.5];
\draw[dashed] (-2.1,0)--(2.1,0) node[right]{$T_z\Sigma$};
\draw[blue!55!black,thick] (-1.15,.13)..controls(-.55,.025)and(-.2,0)..(0,0)
 ..controls(.2,0)and(.55,.025)..(1.15,.13);
\draw[->,thick] (0,0)--(0,.85) node[left]{$\nu(z)$};
\draw[->,thick] (0,0)--(1.5,1.5) node[midway,below right]{$\ell e=y-z$};
\draw[dashed,blue!55!black] (0,1.5)--(1.5,1.5)
 node[midway,above]{$\ell/(2s)$};
\fill (0,0) circle(1.6pt) node[below left]{$z$};
\fill (1.5,1.5) circle(1.6pt) node[right]{$y\in\Sigma$};
\fill[blue!55!black] (0,1.5) circle(1.2pt);
\node[anchor=east] at(-.08,1.65){center};
\end{tikzpicture}
\caption{The normal-centered circle through two zeros, shown on the
side where $\nu(z)\cdot e=s>0$. Its radius is $\ell/(2s)$, so the
unmodified quotient $\ell/s$ is twice that radius. The blue arc shows
only the local sheet at $z$; the picture asserts neither tangency at
$y$ nor that the circle bounds a zero-free ball. The penalized reach
replaces $1/s$ by $\Gamma(s)$ below.}
\label{fig:reach-tangent-ball}
\end{figure}

We slightly modify this quotient when the chord is almost normal.
The angular factor $\Gamma(s)$ will equal $1/s$ away from $s=1$,
stay uniformly close to $1/s$, and be constant near $1$.
Its constant portion makes the second contact normal to the chord;
this is what will permit the signed transfer in
Proposition~\ref{prop:reach-calculus}.  The transition must also have
small weighted second derivative, so its differentiation costs only
an arbitrarily small error.

Here is a penalty with these properties.  Fix an accuracy
$\eps_1>0$, choose $0<\delta_1<1/10$, then $L>1$, and put
$\delta_2=\delta_1e^{-L}$.  Let $\chi:\R\to[0,1]$
be smooth, nonincreasing, equal to $1$ on $(-\infty,0]$, and equal
to $0$ on $[1,\infty)$.  For $s\ge1-\delta_1$ define
\[
 \Gamma(s)=\frac1{1-\delta_1}
 -\int_{1-\delta_1}^s u^{-2}
   \chi\!\left(\frac{\log(\delta_1/(1-u))}{L}\right)\,\dd u,
\]
and set $\Gamma(s)=s^{-1}$ below $1-\delta_1$.  The integrand is
understood as zero near $u=1$.  This gives a smooth function on
$(0,1]$, constant, with value $G_1$, for $s\ge1-\delta_2$.
Direct differentiation gives
\begin{equation}\label{eq:reach-penalty}
 \begin{gathered}
 s^{-1}\le\Gamma(s)\le(1+C\delta_1)s^{-1},\qquad
 1\le G_1\le\Gamma(s),\\
 -(1-\delta_1)^{-2}\le\Gamma'(s)\le0,\qquad
 0\le(1-s)\Gamma''(s)\le
 \frac{2\delta_1}{(1-\delta_1)^3}
 +\frac{\|\chi'\|_\infty}{L(1-\delta_1)^2}
 \le\eps_1
 \quad(s\ge1-\delta_1).
 \end{gathered}
\end{equation}
The last inequality is achieved by first decreasing $\delta_1$ and
then increasing $L$.  Thus both $G_1-1$ and $\eps_1$ can be prescribed
arbitrarily small.  All choices are made before any large reach or
exhaustion limit.

On $\Sigma_{\rm reg}$ define
\begin{equation}\label{eq:reach-definition}
 d(z)=\inf_{y\in\Sigma\setminus\{z\}}\ell\Gamma(s),\qquad
 \ell=|y-z|,\quad e=\frac{y-z}{\ell},\quad
 s=|\nu(z)\cdot e|.
\end{equation}
A competitor with $s=0$ has value $+\infty$.
An \emph{active contact} is a competitor attaining this infimum, or a
limiting diagonal competitor attaining it as the second point approaches
the base.  The proof of Proposition~\ref{prop:reach-calculus} will show
that one of these alternatives always realizes the infimum.

\begin{lemma}[Elementary chord bounds]\label{lem:reach-small-defect}
The function $d$ is finite and positive on $\Sigma_{\rm reg}$.
At every such point, writing $\lambda_i$ for the eigenvalues of $A$,
\begin{equation}\label{eq:reach-elementary}
 d\max_i|\lambda_i|\le2,
 \qquad |(y-z)\cdot\nu(z)|\le C_\Gamma |y-z|^2/d(z)
 \quad(y\in\Sigma).
\end{equation}
For each $J<\infty$ there is a uniform $\delta(J)>0$ such that
$v(z)=0$ and $D(z)<\delta(J)$ imply $d(z)>J$.
Moreover, at every regular base there are zero-free open balls tangent
at that base on both normal sides, of radius $c_\Gamma d(z)$.
\end{lemma}

\begin{proof}
Local smooth graph representation bounds the quotients
$|\nu(z)\cdot(y-z)|/|y-z|^2$ for nearby zeros; competitors outside a
fixed small ball have a positive lower bound because $\Gamma\ge1$.
This proves positivity.  If every zero lay in the tangent plane at
$z$, phase decay \eqref{eq:phase-decay} would bound the energy by
$C e^{-c\dist(x,T_z\Sigma)}$, giving finite density.  Thus a finite
competitor exists.  A signed unit geodesic tangent to a principal
direction gives
$\ell/s\to2/|\lambda_i|$ whenever $\lambda_i\ne0$, proving the first
inequality.  The second uses the \emph{upper} penalty bound:
$d\le\ell\Gamma(s)\le C_\Gamma\ell/s$.

If $D(z_j)\to0$, the translated, rotated fields $U$ converge smoothly
on every fixed ball to $x_n$, by \eqref{eq:diffuse-local} and the
equation.  In $B_{4J}$ all zeros are therefore one graph whose
curvature and slope tend to zero.  A competitor of value at most
$J$ has $\ell\le J$; on this graph $s\le o(1)\ell$, whereas
$\ell\Gamma(s)\ge\ell/s$.  This is impossible, also for competitors
converging to the diagonal.  The contradiction formulation proves
uniformity of the threshold.  Finally, the second inequality in
\eqref{eq:reach-elementary} excludes zeros from each ball
$B_{c_\Gamma d}(z\pm c_\Gamma d\nu)$ when $c_\Gamma$ is sufficiently
small: a point in either ball has
$|\nu\cdot(y-z)|>|y-z|^2/(2c_\Gamma d)$.
\end{proof}

\begin{remark}[Jointly injective ordinary tubes]\label{rem:reach-ordinary-tube}
The elementary chord inequality also gives the ordinary tube used
in the final packing argument.  If $d(z)>L$, choose
$0<c<\min\{1/(2C_\Gamma),1/2\}$.  For $0<|t|<cL$ and every other zero $y$,
\[
 |z+t\nu(z)-y|^2-t^2
 =|y-z|^2-2t\nu(z)\cdot(y-z)
 \ge(1-2cC_\Gamma)|y-z|^2>0.
\]
Thus $z$ is the unique nearest zero to $z+t\nu(z)$.  Normal tubes
over all bases with $d>L$ are therefore jointly injective in this
range, on both sides.  The curvature bound
$\max|\lambda_i|\le2/d$ implies
$\det(\Id+tA)\ge(1-2c)^m>0$.  Consequently any such bases in
$B_{2R}$ have area at most $C R^{m+1}/L$ when $L\le R$, by integrating
the positive Jacobian over both normal sides and
comparing the resulting tube volume to $|B_{3R}|$.
\end{remark}

\subsection{A bounded density window with delayed monotonicity}

A long reach confines nearby zeros to a thin slab.  We use that
confinement to limit the increase of energy density between successive
radii.  The resulting estimate involves the density at a larger radius,
so its discrete iteration must retain that delay.

\begin{lemma}[A fixed density bound near the reach scale]\label{lem:reach-density-window}
There are fixed $C,c>0$ and a fixed density bound $K_0<\infty$
such that, if $d=d(z)$ is sufficiently large, then
\begin{equation}\label{eq:C1}
 h=d\exp\{-C\sqrt{\log(K(d)+1)}\},\qquad M_h(z)\le K_0.
\end{equation}
The constant $C$ can be increased without invalidating the conclusion.
In particular $h\ge d^{1-o(1)}$, uniformly as $d\to\infty$.
\end{lemma}

\begin{proof}
\emph{Slab confinement and one-step density control.}
Translate $z$ to $0$ and take $\nu(z)=e_n$.  By
\eqref{eq:reach-elementary}, the zeros in $B_{3R}$ lie in
$|x_n|\le C R^2/d$.  With $\mathbb S=e(v)\Id-\nabla v\otimes\nabla v$,
the stress tensor identity $\operatorname{div}\mathbb S=0$, tested
against $x_ne_n\zeta^2$, gives
\[
 \int\zeta^2\mathbb S_{nn}
 =-2\int x_n\zeta\,\mathbb S e_n\cdot\nabla\zeta.
\]
The identities $\nabla v=qp/\sqrt2$ and
$e(v)=q^2(1+|p|^2)/4$ show
$\mathbb S_{nn}\ge c q^2(1-p_n^2)$ and
$|\mathbb S e_n|^2\le Cq^2\mathbb S_{nn}$.
Young's inequality, with a cutoff equal to one on $B_R$ and
supported on $B_{2R}$, consequently proves
\[
 \int_{B_R}q^2(1-p_n^2)
 \le C\int_{B_{2R}}q^2(x_n/R)^2.
\]
On the part of $B_{2R}$ with
$|x_n|>CR^2/d+\sqrt R$, every zero either lies outside $B_{3R}$
or has the confined height, so its distance is at least $c\sqrt R$.
Phase decay controls the energy there by $C R^n e^{-c\sqrt R}$.
Elsewhere $(x_n/R)^2\le C((R/d)^2+R^{-1})$.

In the monotonicity formula between $R/B$ and $R$, with a fixed
$B\ge8$, the discrepancy is bounded by $Cq^2(1-p_n^2)$, and the
squared radial derivative on these annuli by
$C_Bq^2(1-p_n^2+(x_n/R)^2)$.  The preceding estimate, at a fixed
multiple of $R$, therefore gives
\begin{equation}\label{eq:C3}
 0\le M_R(z)-M_{R/B}(z)
 \le C\bigl((R/d)^2+R^{-1}\bigr)M_{BR}(z)
 \qquad(10\le R\le c d).
\end{equation}
The exponentially small remainder was absorbed using $M_R(z)\ge c$
for $R\ge1$.  This uniform lower bound follows from $v(z)=0$,
the uniform first derivative bound, and monotonicity.

\emph{Iterating with the larger-radius term.}
We now retain the delay in \eqref{eq:C3}; replacing $M_{BR}$ by
$M_R$ would give an unjustified recursion.
Set $R_j=dB^{-j}$, $k_j=M_{R_j}(z)$, and
$b_j=C(B^{-2j}+R_j^{-1})$.  Beyond a fixed starting index,
\[
 k_j-k_{j+1}\le b_jk_{j-1}.
\]
Choose $R_{\rm fix}$ and that starting index so $b_j<1/8$ while
$R_j\ge R_{\rm fix}$.  Put
$J=\lceil C_0\sqrt{\log(K(d)+1)}\rceil$.
For sufficiently large $d$, $R_{2J+10}\ge d^{1/2}$, because
$K(d)\le1+Cd$.  Suppose a halving $k_{j+1}<k_j/2$ occurs at some
$j\ge J-2$ before the fixed-scale endpoint.  Then
$k_{j-1}>k_j/(2b_j)>4k_j$, so a halving also occurs at the preceding
index.  Induction propagates this backwards through all indices in
$[\lceil J/2\rceil,J-3]$.  On this interval
$b_i\le C e^{-cJ}$: the $B^{-2i}$ term has this bound, and
$R_i^{-1}\le d^{-1/2}$ is smaller.  Multiplying the preceding lower
growth factors over at least $cJ$ indices gives
$K(d)\ge c\exp(c'J^2)$, contradicting the choice of large $C_0$.
Here all $k_j$ before the endpoint are bounded below by the same
positive constant.  Thus no such halving occurs.

It follows that $k_{j-1}\le2k_j$ for $j\ge J$ up to that endpoint,
and hence $k_{j+1}\ge(1-2b_j)k_j$.  The sum of $b_j$ on this
range is uniformly bounded: the first terms form a geometric series,
and the $R_j^{-1}$ terms form a reversed geometric series ending at
a fixed radius.  The product of $(1-2b_j)^{-1}$ is therefore bounded.
At the last index $M_{R_j}$ has a uniform bound from the bounded
energy integrand.  Consequently $k_J\le K_0$ with fixed $K_0$.
Taking $C$ large enough in \eqref{eq:C1} makes $h\le R_J$;
monotonicity completes the proof.
\end{proof}

\subsection{A bounded positive barrier for the normal operator}

The profile error on either side of an isolated zero sheet vanishes
on that sheet.  To compare this error with its forcing we need a
positive bounded supersolution for the normal operator.  The following
formula will also apply when the boundary of a later comparison region
is shifted by a fixed amount.

\begin{lemma}[Positive half-line barrier]\label{lem:reach-half-line-barrier}
Let $g(t)=\tanh(t/\sqrt2)$ and $V_g(t)=3g(t)^2-1$.
For each fixed $a_0\in\R$, the function
\begin{equation}\label{eq:reach-half-line-barrier}
 v_{a_0}(t)=g'(t)\left[1+\int_{a_0}^t(g'(s))^{-2}
                         \int_s^\infty g'(r)\,\dd r\,\dd s\right],
 \qquad t\ge a_0,
\end{equation}
satisfies
\[
 (-\partial_t^2+V_g)v_{a_0}=1,\qquad
 0<c_{a_0}\le v_{a_0}\le C_{a_0}.
\]
Every fixed derivative of $v_{a_0}$ is bounded on this half-line.
The constants can be chosen uniformly when $a_0$ varies over a fixed
compact interval.  Translating the independent variable gives the
same assertion for a translated profile and half-line.
\end{lemma}

\begin{proof}
Write $v_{a_0}=g'b$ and use
$(-\partial_t^2+V_g)g'=0$.  Then
\[
 (-\partial_t^2+V_g)(g'b)
 =-\frac{1}{g'}\bigl((g')^2b'\bigr)'=1,
\]
since $(g')^2b'=\int_t^\infty g'(r)\,\dd r$.
Both terms in the bracket in \eqref{eq:reach-half-line-barrier} are
nonnegative, and $g'>0$.  Moreover
$g'(t)\sim2\sqrt2e^{-\sqrt2t}$ and
$\int_t^\infty g'=1-g(t)\sim2e^{-\sqrt2t}$, so
$v_{a_0}(t)\to1/2$ as $t\to\infty$.
Continuity gives the positive lower and finite upper bounds, uniformly
for $a_0$ in a compact interval.  Differentiating the displayed formula
gives a bounded first derivative; the equation and its successive
derivatives give all remaining fixed derivative bounds.
\end{proof}

\subsection{Isolation and the two parity projections}

The density window now gives one graph on the scale $h$.
For reach calculus we need more than its curvature bound: the first
two derivatives of $H$ must be smaller than the scale-invariant
$h^{-3}$ size.  Two projections in the normal variable provide this
gain.  The first uses oddness of the leading curvature forcing; the
second uses the tangential improvement obtained from the first.

\begin{lemma}[An isolated graph and mean-curvature derivatives]\label{lem:reach-isolation}
There are uniform $c,c',C,\gamma>0$ such that, for sufficiently
large $d(z)$ and $h$ in \eqref{eq:C1}, the zero set in $B_{ch}(z)$
is a single graph over $T_z\Sigma$, extending over the tangent disk
of radius $c'h$.  On a smaller concentric portion,
\begin{equation}\label{eq:C2}
 \begin{gathered}
 |A|+h^{1/2}[A]_{C^{1/2}}\le C h^{-1},\qquad
 |\nabla_\Sigma A|\le C h^{-2},\qquad |H|\le C h^{-2},\\
 |\nabla_\Sigma H|+|\nabla_\Sigma^2H|\le C h^{-3-\gamma}.
 \end{gathered}
\end{equation}
The estimates hold uniformly on every inner patch having a fixed
fractional margin from the original chart boundary.  Hölder norms
may be computed in its graph coordinates.
\end{lemma}

\begin{proof}
\emph{One graph and preliminary bounds.}
The fixed bound of Lemma~\ref{lem:reach-density-window} is subcritical
by Theorem~\ref{thm:finite}.  Apply
Corollary~\ref{cor:finite-subcritical}, including its sheeting conclusions,
at radius $h$.  Thus the density used here is the fixed local bound
$K_0$; the contradiction hypothesis still allows unbounded density on
larger scales.
The confinement thickness divided by this radius is $O(h/d)$.
Increasing the fixed constant in \eqref{eq:C1} makes this ratio as
small as required, without changing the density bound.  The resulting
graphs have arbitrarily small fixed slope, the first estimate in
\eqref{eq:C2}, and
\begin{equation}\label{eq:reach-preliminary-H}
 |H|\le C_\iota h^{-2+\iota}\quad\text{for every fixed }\iota>0.
\end{equation}
An additional graph in an inner cylinder would cross the normal line
at $z$ at distance $O(h)$, producing a competitor of value $O(h)$,
contrary to $d/h$ being sufficiently large.  Thus just one sheet
remains in a smaller cylinder.  On fixed windows centered on inner
zeros the solution, in oriented normal coordinates, converges
uniformly and smoothly to $g$.  The lower gradient bound from sheeting
and interior elliptic estimates also give unit-scale bounds on all
fixed orders of graph derivatives.  These statements are uniform in
the centers and solutions, by the stated uniformity of the subcritical
ball input and the usual sequential formulation of compactness.

It remains to strengthen the mean-curvature derivatives and then recover
the large-scale bound on $\nabla_\Sigma A$.  The two projections below
serve different purposes: the first removes the leading odd forcing
from the equation at the zero sheet, and the second uses the resulting
tangential gain to improve the derivatives of $H$.
We prove the derivative improvement; in particular we do not
differentiate \eqref{eq:reach-preliminary-H} at its original scale.
Fix an inner zero as origin.  Let $X(y)$ be its tangent-plane graph
and use Fermi coordinates $X(y)+t\nu(y)$.  Choose a small number
$\epsilon>0$, to be fixed at the end, and write $R_h=h^\epsilon$.
All arguments below take place in a fixed finite sequence of nested
boxes with side lengths fixed fractions of $R_h$, starting inside
$\{|y|_\infty,|t|<R_h\}$.  The much larger $h$-scale chart guarantees
that these coordinates are injective there, have polynomially bounded
unit-scale derivatives of each fixed order, and contain no other
zero.  The distance to $\Sigma$ is comparable to $|t|$ on the boxes.
Consequently
\[
 \varphi(y,t)=v(X(y)+t\nu(y))-g(t)=o(1)
\]
uniformly.  To see uniformity also at unbounded $t$, first choose a
large fixed depth and use phase decay beyond it, then use the fixed
window convergence at shallower depths.  For $|t|$ a fixed fraction
of $R_h$, $\varphi$ and every needed fixed derivative are bounded by
$C h^C e^{-cR_h}$: apply phase decay and unit-scale elliptic estimates
to $v\mp1$, and differentiate the coordinate composition.

We record the interpolation used below.  If a tensor has uniformly
bounded $C^N$ norm on unit balls and $C^{1/2}$ seminorm at most
$Ch^{-3/2}$, its derivatives of orders $1\le k\le30$ are bounded by
$C h^{-5/4}$, by taking $N$ sufficiently large once and for all.
Indeed the standard interpolation exponent of the low seminorm can
be taken arbitrarily close to one for any fixed $k$ by increasing
$N$; subtracting its value at the center removes the constant term.
Applied to the coordinate components of $A$, whose change-of-frame
terms have the same or better bound, this gives
\begin{equation}\label{eq:reach-interpolation}
 |\partial_y^k A|\le C h^{-5/4}\quad(1\le k\le30),\qquad
 |\partial_y^kH|\le C_\epsilon h^{-2+\epsilon}\quad(0\le k\le30).
\end{equation}
For the second assertion use \eqref{eq:reach-preliminary-H} with
$\iota<\epsilon/2$ and interpolate its supremum norm with a
sufficiently high unit-scale derivative bound.  All interpolation
orders are finite and fixed before $h$ tends to infinity.

In these coordinates the equation is exactly
\begin{equation}\label{eq:C4}
 \begin{split}
 (-\Delta_y-\partial_t^2+V_g)\varphi={}&K_tg'+K_t\varphi_t
 +(g_t^{ij}-\delta^{ij})\partial_{ij}\varphi
 +b_t^i\partial_i\varphi-3g\varphi^2-\varphi^3,\\
 V_g={}&3g^2-1,\qquad
 K_t=\tr[A(\Id+tA)^{-1}]\,.
 \end{split}
\end{equation}
Here $g_t^{ij}$ is the inverse tangential metric, and $b_t^i$ is the
first-order coefficient in its Laplace operator.  Neither superscript
nor subscript $t$ denotes differentiation in this display.  Matrix
inversion in $\Id+tA$, the graph metric, and
\eqref{eq:reach-interpolation} give, with derivatives through order
20,
\begin{equation}\label{eq:reach-coefficients}
 \begin{gathered}
 K_t=H-t|A|^2+t^2\tr A^3+O(h^{-4+C\epsilon}),\qquad
 K_t=O(h^{-2+C\epsilon}),\\
 g_t^{ij}-\delta^{ij}=O(h^{-1+C\epsilon}),\qquad
 b_t^i=O(h^{-1+C\epsilon}).
 \end{gathered}
\end{equation}
Every occurrence of $C\epsilon$ in this proof has a bounded absolute
multiplier independent of $\epsilon$.  For example the tangent slope
is $O(h^{-1+\epsilon})$, and the remainder in the first line contains
at least four factors of $A$ and at most a fixed power of $|t|$.
The differentiated remainders obey the same bounds, since each
positive derivative of $A$ is no larger than $O(h^{-1})$.

\emph{A first bound for the profile error.}
Use $v_0$ from Lemma~\ref{lem:reach-half-line-barrier} on $t\ge0$
and its reflection on $t\le0$, separately on the two half-boxes.
The reflected pair need not be differentiable at zero; no full-line
supersolution is asserted.  In each half-box move the metric
and drift terms to the left of
\eqref{eq:C4} and write the nonlinearity as
$(3g\varphi+\varphi^2)\varphi$.  Its coefficient tends uniformly to
zero, so the resulting operator applied to $v_0$ is at least $1/2$.
Comparison is valid after division by this positive supersolution,
which makes the zeroth-order coefficient positive.  For each lateral
or distant normal face, $v_0e^{-c_0\,\text{distance to face}}$ is also
a supersolution when $c_0$ is a sufficiently small fixed number.
The sum of these face barriers and $C h^{-2+C\epsilon}v_0$ dominates
both signs of $\varphi$.  The bottom boundary has $\varphi(y,0)=0$.
Shrinking the box makes all face contributions exponentially small.
Interior elliptic estimates across $t=0$, using the smooth equation,
then give
\begin{equation}\label{eq:reach-phi-first}
 \|\varphi\|_{C^{15}}\le C_\epsilon h^{-2+C\epsilon}.
\end{equation}
The coefficient norms needed for this bootstrap are uniform by
\eqref{eq:reach-coefficients}; the nonlinearity is smooth and has
already been bounded in supremum norm.

\emph{The inverse on the complement of the translation mode.}
The positive barrier controlled the full profile error.  The next
estimate separates its translation mode from the remaining normal
modes, where a spectral gap is available.
Let $L_t=-\partial_t^2+V_g$.  It has kernel spanned by $g'$ and
a fixed spectral gap $\lambda>0$ on its $L^2(\R)$ orthogonal
complement.  One elementary justification is the ground-state
factorization of its quadratic form, simplicity of the positive
zero mode, and compactness on bounded intervals together with
$V_g\to2$ at infinity; a Rayleigh sequence tending to zero could
neither escape to infinity nor stay orthogonal to the kernel.
Let $P_\perp$ be this projection.  For
\[
 (-\Delta_y+L_t)w=P_\perp F,\qquad P_\perp w=w,
\]
the regularized $L_t^2$ norm of each $y$-derivative of $w$ is a
subsolution for $-\Delta_y+\lambda$, with source the corresponding
norm of $F$.  This follows by differentiating the squared norm,
using $\langle L_tw,w\rangle\ge\lambda\|w\|_2^2$ and
Cauchy--Schwarz, and then letting the regularization vanish.
Comparison with the constant source bound plus decaying face
barriers controls these norms on each smaller $y$-box, with an
exponentially small multiple of their boundary bound.  Commuting two
additional $y$-derivatives and using the equation gives the same
bound in $H_t^2$, since
$\|w\|_{H^2}\le C(\|L_tw\|_2+\|w\|_2)$.
Thus source derivatives through order $j+2$ control the value at
$t=0$ and its $y$-derivatives through order $j$.
This assertion applies separately to even and odd functions of $t$,
because both the operator and $P_\perp$ preserve parity.  It applies
also to a range of strictly positive derivative orders: no estimate
of the undifferentiated source is needed to estimate a differentiated
equation.  Polynomial boundary bounds suffice in every use here.

\emph{The whole-line equation and the first parity projection.}
Choose an even cutoff $\chi_h(t)$, equal to one on an inner normal
interval and supported in a slightly larger interval inside the box,
with derivatives supported where $|t|$ is comparable to $R_h$.  Put
$\Phi=\chi_h\varphi$.  If $E_{\rm raw}$ denotes the right side of
\eqref{eq:C4} minus $Hg'$, the exact equation on the whole normal line is
\[
 (-\Delta_y+L_t)\Phi=Hg'+E_{\rm full},\qquad
 E_{\rm full}=\chi_h E_{\rm raw}+(\chi_h-1)Hg'
                  -2\chi_h'\varphi_t-\chi_h''\varphi.
\]
The kernel tail $(\chi_h-1)Hg'$ and the two commutators, together with
all their needed $y$-derivatives, are $O(h^{-N})$ for every fixed $N$.
For the tail this follows from the preliminary derivatives of $H$ and
the exponential decay of $g'$; for the commutators it follows from the
normal-edge estimates.  Restoring the cutoff leading term
$-\chi_h t|A|^2g'$ to the whole-line odd function $-t|A|^2g'$ costs
another error of the same order.  We henceforth write $\varphi$ for
$\Phi$, and $E$ for $E_{\rm full}$.  Write
\[
 \varphi=u_0(y)g'(t)+\varphi^\perp(y,t),\qquad
 (-\Delta_y+L_t)\varphi=H(y)g'(t)+E(y,t).
\]
We normalize kernel coefficients by $\int(g')^2$.
From \eqref{eq:reach-coefficients} and \eqref{eq:reach-phi-first},
in $L_t^2$ with $y$-derivatives through order 11,
\begin{equation}\label{eq:reach-first-parity}
 E=-t|A|^2g'+O(h^{-3+C\epsilon}),\qquad
 E^{\rm even}=O(h^{-3+C\epsilon}).
\end{equation}
Here $t^2\tr A^3g'$ has size $h^{-3+C\epsilon}$, every metric
error is a coefficient of size $h^{-1+C\epsilon}$ times a derivative
of $\varphi$ of size $h^{-2+C\epsilon}$, and the normal-drift and
nonlinear errors have size $h^{-4+C\epsilon}$.

The inverse estimate for the even orthogonal component, followed by
evaluation at $t=0$, bounds $\varphi^\perp(y,0)$ and its derivatives
through order 9 by $C h^{-3+C\epsilon}$; the odd component vanishes
there.  Since $\varphi(y,0)=0$, the same bound holds for $u_0$.
The kernel equation
\begin{equation}\label{eq:reach-kernel}
 -\Delta_yu_0=H+E_0
\end{equation}
therefore bounds $H$ and its derivatives through order 7 by
$C h^{-3+C\epsilon}$.  In addition, positive $y$-derivatives of
$-t|A|^2g'$ gain $h^{-1/4}$ over its $h^{-2}$ bound by
\eqref{eq:reach-interpolation}; derivatives of the remainder in
\eqref{eq:reach-first-parity} are already smaller.  Applying the
inverse estimate to the entire orthogonal component consequently gives
\begin{equation}\label{eq:reach-positive-phi}
 \|\partial_y^k\varphi\|_{L_t^2}
 \le C h^{-2-1/4+C\epsilon},\qquad1\le k\le8.
\end{equation}
For these $L_t^2$ bounds no loss of two derivative orders is required;
the already controlled derivatives of $u_0$ cover the kernel term.

The first projection has now bounded the kernel coefficient through
nine tangential derivatives and $H$ through seven.  The full orthogonal
equation has also supplied the gain in \eqref{eq:reach-positive-phi}
without losing derivative orders.  We use that gain in a second even
projection; no eighth derivative of $H$ is required, because $P_\perp$
annihilates $Hg'$.

\emph{The second parity projection.}
Return to \eqref{eq:C4}.  Positive derivatives through order 6 of
the even part of $E$ now obey
\begin{equation}\label{eq:reach-second-parity}
 \|\partial_y^kE^{\rm even}\|_{L_t^2}
 \le C h^{-3-1/4+C\epsilon},\qquad1\le k\le6.
\end{equation}
Indeed the odd term $-t|A|^2g'$ still contributes nothing; a positive
derivative of $\tr A^3$ contains a factor bounded by $h^{-5/4}$;
every horizontal metric or drift error contains at least one positive
$y$-derivative of $\varphi$, so \eqref{eq:reach-positive-phi} applies
up to order 8.  If derivatives instead fall on its coefficient, that
coefficient remains $O(h^{-1+C\epsilon})$ and the same estimate
applies to the positive derivative already present on $\varphi$.
The normal drift and nonlinear terms remain $O(h^{-4+C\epsilon})$
in these norms.  This accounts for all terms, including products of
derivatives and the superpolynomially small cutoff errors.

Apply the even orthogonal inverse once more.  Its value at zero and
the zero condition give
$|\partial_y^ku_0|\le C h^{-3-1/4+C\epsilon}$ for $1\le k\le4$.
Equation~\eqref{eq:reach-kernel} gives the same estimate for
$\partial_yH$ and $\partial_y^2H$.  Choose $\epsilon$ so small that
every fixed exponent loss $C\epsilon$ above is less than $1/8$;
decreasing the resulting gain if necessary gives $\gamma>0$.
Covariant derivatives have the asserted bounds as well: the coordinate
connection is bounded and its extra terms multiply already controlled
first derivatives of $H$.

\emph{Recovering the large-scale curvature derivative.}
The improved mean-curvature derivatives are now in hand.  Rescale
the original graph disk by $h^{-1}$.  Its existing
$C^{2,1/2}$ bound is uniform, and the rescaled mean curvature $hH$
has uniformly bounded $C^{1,1}$ norm, since its first and second
derivatives scale by $h^2$ and $h^3$, respectively.  Differentiate
the prescribed mean-curvature graph equation once and apply interior
Schauder estimates.  The rescaled graph has a uniform third derivative
bound, giving $|\nabla_\Sigma A|\le Ch^{-2}$ in original variables.
All boxes used above have a fixed total fractional shrinkage and all
derivative orders were fixed in advance.  Repeating at each inner
center proves the uniform inner-window assertion.
\end{proof}
\subsection{One-sided trace improvement from a tangent phase ball}

The isolated chart controls the base point of a reach contact.  Its
second point need not lie in that chart, so it requires an estimate
whose sole geometric input is a zero-free tangent ball.  The sign of
the resulting trace bound matters: only positive mean curvature at the
receiving point will enter the upper reach inequality.

\begin{lemma}[One-sided improvement at a tangent phase ball]\label{lem:reach-phase-ball}
There are uniform $r_0,C,\gamma>0$ with the following property.
If an open ball of radius $r\ge r_0$ is disjoint from $\Sigma$ and
its boundary contains $x\in\Sigma$, and $N$ is its inward normal at
$x$, then
\begin{equation}\label{eq:C5}
 x\in\Sigma_{\rm reg},\qquad D(x)\le C/r,
 \qquad a(x)\le C/r,\qquad H_N(x)\le Cr^{-1-\gamma}.
\end{equation}
After changing the sign of $v$ so that it is positive inside the
ball, along the inward radius one also has
\begin{equation}\label{eq:C6}
 U(x+sN)\ge s-C(1+s)/r\qquad(0\le s\le\log r).
\end{equation}
Only the last curvature estimate is one-sided; no absolute-value
improvement for $H_N$ is asserted.
\end{lemma}

\begin{proof}
\emph{Comparison along the inward radius.}
Let $s$ denote inward depth in the annulus $0<s<Q=\sqrt r$.
The nonnegative function
\[
 w(s)=g'(s)\int_0^s(g'(t))^{-2}
                    \int_t^\infty(g'(u))^2\,\dd u\,\dd t
\]
has $w(0)=0$, $(-\partial_s^2+V_g)w=g'$, and
$|w|+|w'|\le C(1+s)g'$.  Thus $v_*=g-C_0w/r$, with a sufficiently
large fixed $C_0$, is positive on the annulus and satisfies
$\Delta v_*+v_*-v_*^3>0$.  Indeed the linear correction contributes
$C_0g'/r$, the inward radial drift is
$-(n-1)v_*'/(r-s)$, and the quadratic correction is bounded by
$C r^{-2}(1+s)^2(g')^2=o(r^{-1}g')$ uniformly on $[0,Q]$.
Multiplication by $1-e^{-cQ}$, with a sufficiently small fixed
$c>0$, keeps a strict subsolution and makes its value less than $v$
on the inner boundary, by phase decay.  On the outer boundary it
vanishes.  Sliding its amplitude from a small positive number to one
proves comparison.  To justify the start, $v>0$ in the compact
annulus and the Hopf boundary estimate applies at any outer zero;
these and compactness bound the quotient from below.  At a first
contact a strict subsolution cannot touch in the interior, and the
Hopf comparison excludes an outer boundary obstruction.  Smaller
amplitudes remain strict subsolutions: writing
$\mathcal F(w)=\Delta w+w-w^3$, one has
$\mathcal F(av_*)=a\mathcal F(v_*)+(a-a^3)v_*^3>0$ for $0<a\le1$.

Inverting $g$ on $0\le s\le\log r$ proves \eqref{eq:C6}:
the correction is bounded by $C(1+s)g'(s)/r$, and the extra
$e^{-c\sqrt r}$ term is smaller after division by $g'(s)$ on that
range.  Differentiation at the contact point gives
$\partial_Nv(x)\ge g'(0)-C/r>0$.  Thus $x$ is regular, $N=\nu$
after the sign choice, and $D(x)\le C/r$.  The tangent ball gives
the principal-curvature lower bound $\lambda_i(A_N)\ge-1/r$.
The zero-level identity $H_N=\partial_ND/(2|p|^2)$ and
\eqref{eq:diffuse-local} give $|H_N|\le CD\le C/r$.
Summing the principal curvatures then bounds each one from above by
$C/r$, proving the full curvature estimate in \eqref{eq:C5}.

\emph{Separating the even source from the boundary contribution.}
The comparison has established regularity and the $r^{-1}$ bounds.
To improve only the upper trace, place $x=0$, $N=e_n$, and use the
ambient vertical coordinate $s$.  Work in $B_P$ with
$P=\alpha\log r$, where $\alpha>0$ is fixed sufficiently small at
the end.  The differentiated diffuse estimates give, for every
needed fixed order $k$,
\[
 \|D\|_{C^k(B_P)}\le C_k r^{-1+C\alpha},\qquad
 \|T\|_{C^k(B_P)}+\|U-s\|_{C^k(B_P)}
 \le C_k r^{-1/2+C\alpha}.
\]
These follow by applying unit-ball derivative estimates and the
exponential comparison for $D$ along chains from $0$; integration
of $T$ gives the estimate on $U-s$.  The multipliers of $\alpha$
are fixed independently of small $\alpha$.  The identity
$\nabla_pT=-T^2-\nabla^2D/2$ further gives
$\partial_sT=O(r^{-1+C\alpha})$: the replacement of $p$ by $e_n$
has error $|p-e_n||\nabla T|=O(r^{-1+C\alpha})$.

Set $F=G(s)D$.  Conjugating the defect equation and replacing $vp$
by $g(s)e_n$ and $q|p|^2$ by $G(s)$ yields
\begin{equation}\label{eq:reach-ball-conjugation}
 (-\Delta+2-G(s))F
 =2G(s)|T(x',0)|^2+\mathcal E,
 \qquad\|\mathcal E\|_\infty\le Cr^{-3/2+C\alpha}.
\end{equation}
For completeness, each replacement in the drift or potential
multiplies a coefficient of size $r^{-1/2+C\alpha}$ by $D$ or
$\nabla D$ of size $r^{-1+C\alpha}$.  Replacing
$|T(x',s)|^2$ by $|T(x',0)|^2$ has the same order, since
$|T|=O(r^{-1/2+C\alpha})$, $\partial_sT=O(r^{-1+C\alpha})$,
and $|s|\le P$.  Logarithmic factors are absorbed by increasing
the fixed multiplier of $\alpha$.

Solve the Dirichlet problem on $B_P$ by linearity.  The solution
with source $2G|T(x',0)|^2$ and zero boundary is even in $s$ and
nonnegative.  The zero-boundary solution for $\mathcal E$ has
interior $C^1$ norm at most $Cr^{-3/2+C\alpha}$, since the potential
$2-G$ is at least one.  The remaining homogeneous solution is
$u_++u_-$, both nonnegative, with boundary values $F$ on the upper
and lower hemispheres, respectively, and zero on the other one.
Their harmless equatorial boundary discontinuity is interpreted by
bounded Dirichlet solutions; it does not affect interior estimates.
Reflection and the maximum principle in the upper half-ball give
$\partial_su_-(0)\le0$; local Harnack and derivative estimates give
$|\partial_su_+(0)|\le Cu_+(0)$.

At $x_*=(0,P/2)$, positivity of the source solution implies
\[
 u_+(x_*)\le F(x_*)+Cr^{-3/2+C\alpha}
 \le C(1+P)r^{-1}G(P/2)+Cr^{-3/2+C\alpha}.
\]
For the last inequality integrate $D\le2(1-\partial_sU)$ along
$[0,P/2+1]$, use \eqref{eq:C6}, and apply unit-scale Harnack to $D$
to pass from the integral to its value at $P/2$.  This use is valid
for $\alpha<1$, so the interval lies in the range of \eqref{eq:C6}.

\emph{Suppressing the upper boundary contribution.}
Only $u_+$ can contribute a positive leading normal derivative at the
origin.  Its value at $x_*$ is small by the inward comparison; we
now transfer that gain to the origin with explicit Poisson bounds.
Write $A_\partial=\int_{\partial B_P\cap\{s>0\}}F\,\dd S$.
Comparison with potentials one and two gives
\begin{equation}\label{eq:reach-poisson}
 u_+(0)\le CP^C e^{-P}A_\partial,
 \qquad
 u_+(x_*)\ge c_\zeta
 e^{-(\sqrt2+\zeta)\sqrt5 P/2}A_\partial
 \quad(\zeta>0).
\end{equation}
For the first inequality, the Poisson density at the center for
$-\Delta+1$ is constant on the sphere.  The spherical average of
$e^{x\cdot\theta}$ shows that its size is at most $CP^Ce^{-P}$.
The solution for potential $2-G\ge1$ is smaller.  For the second,
the solution for potential $2-G\le2$ is larger than that for
$-\Delta+2$.  Its Dirichlet Poisson density at any upper-hemisphere
point $\xi$ is at least
$c_\zeta e^{-(\sqrt2+\zeta)|x_*-\xi|}$.  Here are quantitative
details of this last bound.  Start the Dirichlet Green function a
fixed distance from $x_*$, with a fixed positive lower bound obtained
from a fixed-radius ball.  Along the segment towards $\xi$ insert
a cylinder of fixed large cross-sectional radius, chosen so that
$\sqrt{2+\lambda_1}<\sqrt2+\zeta$.  Its first transverse
Dirichlet mode times the longitudinal hyperbolic-sine solution
propagates the lower bound with that exponential rate.  Stop a fixed
distance before $\xi$.  Such a cylinder fits in $B_P$ except for
fixed initial and final portions: $x_*$ is at distance $P/2$ from
the boundary and the incidence angle at $\xi$ has a cosine bounded
below uniformly.  Harnack on those fixed portions and a tangent-ball
barrier at $\xi$ transfer the bound to the inward normal derivative.
All constants depend on $\zeta$ but not $P$.  Finally
$|x_*-\xi|\le\sqrt5 P/2$, proving \eqref{eq:reach-poisson}.

Combining the two bounds with $G(P/2)\le Ce^{-\sqrt2 P/2}$ gives
\[
 u_+(0)\le CP^C r^{-1}
 e^{[(\sqrt2+\zeta)\sqrt5/2-1-\sqrt2/2]P}
 +CP^C r^{-3/2+C\alpha}
 e^{[(\sqrt2+\zeta)\sqrt5/2-1]P}.
\]
Since $\sqrt{10}/2-1-\sqrt2/2<0$, choose $\zeta$ so the first
exponent is negative, then choose $\alpha>0$ so the second term
still has a power better than $r^{-1}$.  Polynomial factors in $P$
can be absorbed by slightly decreasing a resulting gain $\gamma>0$.
The source solution has zero derivative at $0$, the lower-hemisphere
solution has nonpositive derivative there, and the error solution is
smaller still.  Thus $F_s(0)\le Cr^{-1-\gamma}$.  Since
$G(0)=1$, $G'(0)=0$, and $|p(0)|^2=1-O(r^{-1})$, the zero-level
identity for $H_N$ proves \eqref{eq:C5}.
\end{proof}

\subsection{The weak reach inequality and transfer of positive charge}

We now combine the base estimates with the tangent-ball bound.  At an
active contact on the plateau of the penalty, the second point will be
called the \emph{receiving foot}.
The proposition distinguishes these contact types because only plateau
contacts retain a mean-curvature term at that foot.  It also specifies
the direction and multiplicity of the area change needed when that
positive term is integrated.

\begin{proposition}[Weak reach inequality and positive-foot transfer]\label{prop:reach-calculus}
Given $\eta_0>0$, the penalty parameters can be chosen, in the order
specified above, so that the following holds for all sufficiently
large $d$.  The function $d$ is locally semiconcave on
$\Sigma_{\rm reg}$, and the singular part of its distributional
surface Laplacian is nonpositive.  Almost everywhere,
\begin{equation}\label{eq:C7}
 |\nabla_\Sigma d|\le
 C_\Gamma\exp\{C_\Gamma\sqrt{\log(K(d)+1)}\}.
\end{equation}
At an active non-plateau or diagonal contact,
\begin{equation}\label{eq:C8}
 \Delta_\Sigma d\le(1+\eta_0+C_\Gamma d^{-\gamma})a^2d
                  +B\,\frac{|\nabla_\Sigma d|^2}{d},
 \qquad B\le1+C_\Gamma d^{-\gamma}.
\end{equation}
If $ad$ is sufficiently small, the same bound can be used with
$B\le\eta_0+C_\Gamma ad+C_\Gamma d^{-\gamma}$, and then
$|\nabla_\Sigma d|\ge c_\Gamma>0$.
At an active plateau contact, meaning $s\ge1-\delta_2$, one has
\begin{equation}\label{eq:C9}
 \Delta_\Sigma d\le
 G_1\bigl(|H(z)|+H_{N_y}(y)_+\bigr)
 +\eta_0\frac{|\nabla_\Sigma d|^2}{d},\qquad N_y=-e.
\end{equation}
For any measurable choice of such contacts, the positive foot charge
can be transferred to $y$ with multiplicity at most one and area factor
\begin{equation}\label{eq:C10}
 J_{y\to z}=
 \frac{|\det(\Id+\ell A_{N_y}(y))|}{s}\le C_\Gamma,
 \qquad d(y)\le d(z)=G_1\ell,
 \qquad a(y)+D(y)\le C_\Gamma/d(z).
\end{equation}
More precisely, $J_{y\to z}\le1+C\delta_2+C a(y)\ell$ whenever
$a(y)\ell$ is bounded by a sufficiently small fixed constant.  The
constants and the large-reach threshold are fixed before any later
defect threshold, collar scale, or exhaustion radius is chosen.
\end{proposition}

\begin{proof}
\emph{Semiconcavity, including diagonal competitors.}
We first justify the weak interpretation, independently of the
differentiations to follow.  On a compact regular coordinate patch
there are uniform bounds $0<c\le d\le C$: local graph curvature
gives the lower bound, and one fixed off-diagonal positive-angle
competitor gives the upper bound after shrinking the patch.
Relevant competitors of value at most $2C$ have bounded length.
Those of length at least $\varepsilon>0$ have angle
$s\ge\varepsilon/(2C)$, by $\Gamma\ge1/s$.  With the second foot
fixed, these are smooth branches in the base variable, uniformly
bounded in $C^2$; no regularity of the second foot is needed here.

For shorter chords the second point is on the same regular local
graph.  In coordinates $X(u)$ write it as $X(u+t\omega)$, with
$|\omega|=1$ and $t$ small.  Such chords use $\Gamma=1/s$.
The reciprocal branch is
\[
 \frac{|\nu(u)\cdot[X(u+t\omega)-X(u)]|}
 {|X(u+t\omega)-X(u)|^2}.
\]
Both numerator before absolute value and denominator vanish to
second order in $t$, since $DX$ is tangent.  Cancelling $t^2$ by
the integral Taylor formula gives a $C^2$ function through $t=0$,
uniformly in the compact direction parameter.  On relevant branches
its absolute value is bounded below by $1/(2C)$, so a smooth sign
can be chosen and taking its reciprocal preserves the uniform
$C^2$ bound.  The diagonal limits are legitimate upper branches
for $d$, since they are limits of competitors.  Compactness of the
remaining parameter sets shows that every minimizing sequence
has either an off-diagonal or a diagonal realizing parameter and
that $d$ is the infimum of a family with locally uniform upper
Hessian bound.  Subtracting the corresponding fixed quadratic
function makes each branch concave and hence also their infimum
concave.  This proves local semiconcavity.

At almost every point there is an Alexandrov second-order expansion.
At such a point any smooth active upper branch has the same first
derivative and an upper bound for the Hessian of $d$.  Its trace
therefore bounds the absolutely continuous Laplacian.  The singular
part is a nonpositive measure by semiconcavity, in coordinates and
therefore for the surface Laplacian as well.  It suffices to compute
the following upper touching branches.

The weak interpretation is now established.  We next compute smooth
upper touching branches, treating the unmodified penalty, diagonal
limits, and the transition to the plateau in that order.  The final
step will transfer the positive plateau contribution to its receiving
foot.

\emph{A tangent phase ball at every second foot.}
For an off-diagonal contact choose $N=\pm\nu(z)$ so $N\cdot e=s>0$.
The open region
\[
 \{z+\ell'e':N\cdot e'>0,\ \ell'\Gamma(N\cdot e')<d\}
\]
contains no zeros.  After translation, rotation, and division by $d$,
its boundary is a fixed smooth embedded compact radial surface.  Near
the origin it agrees with the sphere $|x|^2=x_n$, because there
$\Gamma(s)=1/s$.  On the remainder it is the smooth positive radial
graph $1/\Gamma$.  It has an inward rolling ball of a fixed positive
radius $c_\Gamma$: local $C^2$ graph bounds supply the balls near
each point, and compactness and embeddedness allow a uniform
decrease so those balls avoid the rest of the boundary.  Hence at
the second foot $y$ there is a zero-free tangent ball of radius
$c_\Gamma d$.  Lemma~\ref{lem:reach-phase-ball} shows that the foot
is regular, $a(y)+D(y)\le C_\Gamma/d$, and its inward trace satisfies
$H_{N_y}(y)\le C_\Gamma d^{-1-\gamma}$.

\emph{Off-diagonal contacts with the unmodified penalty.}
First suppose $s\le1-\delta_1$, so $\Gamma=1/s$.  Write
$r=d/2$, $b=e-sN$, and $M=\Id+rA_N$.  Then $\ell=2rs$ and the
inward normal at the second foot is $N_y=N-2se$.
Reflect an orthonormal frame $v_i$ of $T_z\Sigma$ across $e^\perp$
to obtain $v_i'=v_i-2(e\cdot v_i)e$ in $T_y\Sigma$.  Move $z$ and
$y$ on the respective geodesics with these initial velocities.
The resulting chord branches touch $d$ from above.  First and second
differentiation gives
\begin{equation}\label{eq:C11}
 \begin{split}
 \nabla_\Sigma r={}&-Mb/s,\\
 \Delta_\Sigma d\le{}&a^2d+
 \left(2-\frac2{1+r\max_i|\lambda_i|}\right)
             \frac{|\nabla_\Sigma d|^2}{d}\\
 &+\frac{H_{N_y}(y)-(1-2s^2)H_N(z)
                -\ell b\cdot\nabla_\Sigma H_N(z)}{s^2}.
 \end{split}
\end{equation}
Here is the algebra, including the factor of the gradient term.
Put $w=y-z$ and $w_i=v_i'-v_i=-2(e\cdot v_i)e$ and differentiate
$|w|^2/2-rN\cdot w$, which vanishes at contact and is nonnegative
for a competitor.  The summed second derivative bounds
$(\Delta r)s\ell$ by
\[
 \begin{split}
 &\sum_i|w_i|^2+(w-rN)\cdot(-H_{N_y}N_y+H_NN)
 -2r\sum_iw_i\cdot A_Nv_i\\
 &\hspace{8mm}-rw\cdot(\nabla H_N-a^2N)
 -2\sum_i r_i\bigl(w_i\cdot N+w\cdot A_Nv_i\bigr).
 \end{split}
\]
Use $w-rN=-rN_y$ and the first derivative identity.
The algebraic terms besides $a^2$ and the traces reduce to
$4s^2|\nabla r|^2+4s b\cdot\nabla r$.  By
\eqref{eq:reach-elementary}, $0\le M\le
(1+r\max_i|\lambda_i|)\Id$, so
\[
 b\cdot\nabla r=-b\cdot Mb/s
 \le-\frac{s|\nabla r|^2}{1+r\max_i|\lambda_i|}.
\]
Dividing the preceding second derivative expression by $rs^2$
therefore gives exactly \eqref{eq:C11}.  The identity
$\Delta_\Sigma N=\nabla_\Sigma H_N-a^2N$ fixes all trace signs.

We bound its last quotient uniformly, including almost tangential
contacts.  Choose a small fixed $\gamma'>0$ compared with the gains
in Lemmas~\ref{lem:reach-isolation} and \ref{lem:reach-phase-ball}.
If $s\ge d^{-\gamma'}$, the second-foot upper trace, the base bounds
$|H|\le Ch^{-2}$ and $|\nabla H|\le Ch^{-3-\gamma}$, and
$h=d^{1-o(1)}$ bound that quotient above by
$C_\Gamma d^{-1-\gamma''}$ for some fixed $\gamma''>0$.
If $s<d^{-\gamma'}$, then $\ell=ds$ lies inside the isolated graph.
The two normal orientations agree by continuation there.  In its
tangent-plane coordinates the displacement is exactly $\ell b$;
Taylor's formula gives
\[
 H_{N_y}(y)-H_N(z)-\ell b\cdot\nabla H_N(z)
 =O(\ell^2h^{-3-\gamma}).
\]
The coordinate Hessian obeys this bound as well, since its connection
correction multiplies the already controlled gradient.  Dividing by
$s^2$ and adding $2H_N(z)$ proves the same upper error after possibly
decreasing $\gamma''$.

The coefficient in parentheses in \eqref{eq:C11} is at most one,
because $r\max|\lambda_i|\le1$.  It is $O(ad)$ if $ad$ is small.
Also $|b|\ge\sqrt{1-(1-\delta_1)^2}>0$.  If $ad$ is sufficiently
small then the first derivative identity gives
$|\nabla d|\ge c_\Gamma$.  Otherwise $a^2d\ge c_\Gamma/d$.
In either case the error $C d^{-1-\gamma''}$ can be absorbed into
$C d^{-\gamma''}(a^2d+|\nabla d|^2/d)$.  This proves
\eqref{eq:C8} in the present range.

For the gradient bound, if $\ell\ge ch$, then
$1/s=d/\ell\le Cd/h$ and $\|M\|\le2$, giving
$|\nabla d|\le Cd/h$.  For $\ell<ch$ both points lie on the
isolated graph.  Expanding its normal and projecting tangentially
at the base yields
\[
 -2s b=\ell A_Nb+O(\ell^2/h^2),
\]
using $|\nabla A|\le Ch^{-2}$ and $|A|\le Ch^{-1}$.
Thus $|Mb|\le Cs d^2/h^2$, and
$|\nabla d|\le C d^2/h^2$.  Both estimates imply \eqref{eq:C7}.

\emph{Diagonal contacts.}
At a diagonal contact the curvature itself is the active constraint.
Choose $N$ so $S=-A_N$ has
top positive eigenvalue $\lambda=2/d$, and let $v_1$ be an active
top eigenvector.  Every eigenvalue $\lambda_j$ of $S$ is between
$-\lambda$ and $\lambda$.  For a unit vector $v$ in the top
eigenspace $V$, expansion along the geodesic with signed length $L$
gives
\[
 \frac{2N\cdot(y-z)}{|y-z|^2}
 =\lambda+\frac13\nabla_vS(v,v)L+O(L^2).
\]
This reciprocal cannot exceed $2/d=\lambda$ for either sign of
$L$, so $\nabla_vS(v,v)=0$.  Codazzi makes $\nabla S$ totally
symmetric, and polarization implies that it vanishes on triples
in $V$.

Touch $d$ from above by $2/S(v_1,v_1)$ using a unit field extension.
At the point choose its derivatives, in a principal frame, to be
\[
 \langle\nabla_i v_1,v_j\rangle
 =\frac{\nabla_iS_{1j}}{\lambda-\lambda_j}
 \quad(\lambda_j<\lambda),
 \qquad \langle\nabla_i v_1,v_j\rangle=0\quad(v_j\in V).
\]
A smooth local unit extension with these specified first derivatives
exists; its second derivatives are unconstrained apart from the unit
length condition, and their tangential components drop out of the
Laplacian of $S(v_1,v_1)$ at a top eigenvector.  Thus that Laplacian is
\[
 (\Delta S)_{11}
 +2\sum_{\lambda_j<\lambda}
             \frac{|\nabla S_{1j}|^2}{\lambda-\lambda_j}.
\]
The Euclidean hypersurface identity
\begin{equation}\label{eq:reach-simons}
 \Delta S=\nabla^2\tr S+(\tr S)S^2-a^2S
\end{equation}
follows by commuting the two derivatives in Codazzi and applying
the Gauss equation.  This is the Simons curvature calculation
\cite{Simons1968,SchoenSimonYau1975}, with the mean-curvature terms
retained because the sheet need not be minimal.
More explicitly the Ricci contraction is
$((\tr S)S-S^2)_{ip}S_{pj}$ and the remaining curvature contraction
is $S^3_{ij}-a^2S_{ij}$, giving the stated identity.  Also
$\lambda-\lambda_j\le2\lambda$, and Codazzi gives
$\nabla_jS_{11}=\nabla_1S_{j1}$, which is zero for $v_j\in V$.
Consequently the denominator $f=S(v_1,v_1)$ satisfies
\[
 \Delta f\ge-a^2\lambda
 -C(|\nabla^2H|+a^2|H|)+|\nabla f|^2/\lambda.
\]
Differentiating $2/f$ now gives
\[
 \Delta d\le a^2d+\frac{|\nabla d|^2}{d}
 +C d^2(|\nabla^2H|+a^2|H|).
\]
The last term is $C d^{-\gamma}(a^2d)$ by
Lemma~\ref{lem:reach-isolation}, $h=d^{1-o(1)}$, and
$a^2d\ge4/d$.  This proves \eqref{eq:C8}.  A diagonal contact
always has $ad\ge2$, so the small-$ad$ assertion makes no claim here.
Finally $|\nabla d|\le Cd^2|\nabla A|\le Cd^2/h^2$ proves
\eqref{eq:C7} for these contacts too.

\emph{The angular transition and the plateau.}
The unmodified and diagonal cases satisfy the claimed inequality.
It remains to handle $s\ge1-\delta_1$, where the penalty changes
the variation at the second foot.  Set
\[
 b=e-sN,\qquad c_1=\Gamma-s\Gamma',\qquad
 Z=c_1e+\Gamma'N.
\]
First variation in the second foot shows that its inward normal is
$N_y=-Z/|Z|$.  Rotate the tangent frame at $z$ by the small plane
rotation taking $N$ to $Z/|Z|$, to obtain $v_i'$ at $y$.
With $w_i=v_i'-v_i$, elementary plane-rotation formulas give
\[
 w_i=O(|b|)N+O(|b|^2),\qquad
 P_{e^\perp}w_i=O(|b|^2),\qquad
 (N-se)\cdot w_i=O(|b|^3).
\]
The constants are uniform for the fixed penalty; in fact
$|Z|=1+O(\delta_1)$ and the rotation is uniformly small.  At $b=0$
use the identity rotation.  Differentiation of $\ell\Gamma(s)$
gives the first derivative and the following upper trace:
\begin{equation}\label{eq:reach-transition-calculus}
 \begin{split}
 \nabla d={}&(-c_1\Id+\ell\Gamma'A_N)b,\\
 \Delta d\le{}&|Z|H_{N_y}+H_N Z\cdot N
 +\ell\Gamma' b\cdot\nabla H_N-\ell\Gamma's a^2\\
 &+\sum_i\left\{
 \frac{c_1}{\ell}|P_{e^\perp}w_i|^2
 +2\Gamma' w_i\cdot A_Nv_i
 +\ell\Gamma''\left[b\cdot A_Nv_i
       +\frac{(N-se)\cdot w_i}{\ell}\right]^2\right\}.
 \end{split}
\end{equation}
This can also be derived directly from
$\dd\ell=e\cdot\dd w$ and
$\dd s=\dd N\cdot e+(N-se)\cdot\dd w/\ell$:
the Hessian in $w$ contributes the first term in braces,
the two mixed normal derivatives the second, the second derivative
of $\Gamma$ the square, and
$\Delta N=\nabla H_N-a^2N$ gives the two normal-variation terms
on the first line of the trace.  The accelerations of the two
geodesic frames give the two trace terms.

Since $\Gamma''|b|^2\le C\eps_1$, the algebraic terms, including
$-\ell\Gamma'sa^2$, are at most
\[
 (1+C(\eps_1+\delta_1))a^2d
       +Ca|b|^2+C|b|^4/d.
\]
For clarity the leading coefficient is bounded by one up to
$O(\delta_1)$ because $-s\Gamma'/\Gamma\le1+C\delta_1$;
the $\Gamma''$ square contributes $C\eps_1a^2d$ and
$C\eps_1|b|^4/d$, and the rotation gives the remaining two terms.
Young's inequality bounds $Ca|b|^2$ by any fixed small multiple
of $a^2d$ plus a fixed multiple of $|b|^4/d$.
If $ad$ is small, the first derivative formula bounds
$|\nabla d|$ below by a fixed multiple of $|b|$.
If it is not small, $|b|^4/d$ is at most
$C\delta_1^2a^2d$.  Choosing $\eps_1$ and then $\delta_1$ small
enough for the prescribed tolerance therefore bounds all these
terms by $(1+\eta_0)a^2d+\eta_0|\nabla d|^2/d$, with room to
decrease the initial tolerance for the subsequent errors.
In the transition interval before the plateau, $|b|$ has a fixed
positive lower bound depending on $\delta_2$.  The trace terms in
\eqref{eq:reach-transition-calculus} are
$O(d^{-1-\gamma})$ from above by the base and second-foot estimates.
They are absorbed exactly as in \eqref{eq:C11}, proving
\eqref{eq:C8}, including its small-$ad$ statement.  The gradient
formula, $\ell\le d$, and $da\le2\sqrt m$ give a fixed bound in
this range, hence \eqref{eq:C7}.

On the plateau $\Gamma'=\Gamma''=0$, so
$\nabla d=-G_1b$, $Z=G_1e$, $N_y=-e$.
The only algebraic remainder is $C|b|^4/d$, bounded by
$C\delta_2|\nabla d|^2/d$.  The traces are
$G_1H_{N_y}+G_1sH_N\le G_1(H_{N_y,+}+|H_N|)$.
Decreasing $\delta_2$ if necessary proves \eqref{eq:C9}.
At the endpoints of the ranges the jets match, since the constructed
penalty is smooth; either computation thus has the same limiting
validity.

The differential inequalities have now been proved for every active
contact type.  It remains to justify their integrated use: a positive
plateau contribution at a receiving foot must be counted with its
inverse area factor, and without an additional multiplicity.

\emph{The inverse Jacobian and positive receiving multiplicity.}
We prove the transfer statement in the direction from receiving feet
to bases; this direction determines the Jacobian below.
For a plateau pair, $z=y+\ell N_y$ and
$|N_y\cdot\nu(z)|=s\ge1-\delta_2$.  Fix a local regular sheet
through $z$.  The equation that $y+tN_y$ lies on this sheet has
nonzero $t$ derivative, so the implicit function theorem supplies
a smooth inverse parametrization
$z(y)=y+\ell(y)N_y$ near this pair.  Projecting its derivative onto
$T_y\Sigma$ removes $(\dd\ell)N_y$ and leaves
$\Id+\ell A_{N_y}$.  Orthogonal projection from $T_z\Sigma$ to
$T_y\Sigma$ has area factor $s$.  It follows that the area
Jacobian of the inverse map is exactly the quotient in
\eqref{eq:C10}.  The second-foot bound $a(y)\le C_\Gamma/d(z)$
therefore bounds this quotient uniformly.  Expanding the determinant
when $a(y)\ell$ is small gives the more precise estimate stated.
The reverse chord is normal at $y$, so it is a competitor of value
$G_1\ell$ there and proves $d(y)\le d(z)$.

At a foot with $H_{N_y}>0$ the orientation $N_y$ is unique.  All
bases sending positive charge to this foot must lie on its same
normal ray.  Suppose two existed, at lengths $0<\ell_1<\ell_2$.
For the farther base $z_2$, the nearer base $z_1$ is a competing
zero with chord direction exactly equal to the direction towards
$y$ and length $\ell_2-\ell_1$.  It has the same angle penalty and
strictly smaller objective than the alleged active chord to $y$,
a contradiction.  Thus every positively charged foot has at most
one inverse base.

There is no hidden diffeomorphism assumption here.  Cover the regular
pair relation by countably many of the preceding inverse charts.
Active parameters on compact coordinate pieces form a closed
compact-valued relation; a measurable selection can be obtained
there and pieced together over a countable exhaustion.  Restrict the
selected pairs to disjoint measurable pieces in those charts.  The
area formula on each inverse chart gives exactly its Jacobian; its
critical image has zero surface measure.  Summing the formula over
the selected pieces adds no multiplicity for positive charges by
the ray argument.  In explicit terms, for a nonnegative measurable
integrand $F(z,y)$ supported on positively charged selected pairs,
\[
 \int F(z,y(z))\,\dd S_z
 =\int F(z(y),y)J_{y\to z}\,\dd S_y,
\]
where the right-hand integral is over the selected positive feet,
with null critical images omitted on the left.  This proves the
claimed transfer for arbitrary nonnegative weights, not only
constant ones, and completes the proposition.
\end{proof}

\subsection{Defect and reach near bad zeros}

We finish with the collar estimate used when geometric tests are
joined to ambient tests.  Its source-reach bound also applies when
the second contact, rather than the base, lies close to a bad zero.

For the next statement, a \emph{good zero} relative to a set
$X\subset\Sigma$ means a point of $\Sigma_{\rm reg}\setminus X$.

\begin{corollary}[Reach bounds near bad zeros]\label{cor:reach-base-defect}
For sufficiently large $d(z)$, $D(z)+a(z)\le C_\Gamma/d(z)$.
If $X=\{z\in\Sigma:D(z)\ge\delta\}$ with fixed $\delta>0$, then
for every fixed $L_0$ there is a finite uniform upper bound for $d(z)$
at all good zeros with $\dist(z,X)\le10L_0$.  The same conclusion
holds for a contact reach $d(z)$ when its second foot $y$ satisfies
$\dist(y,X)\le10L_0$.
\end{corollary}

\begin{proof}
Use the tangent balls of Lemma~\ref{lem:reach-small-defect} and
Lemma~\ref{lem:reach-phase-ball}.  The proof of
Proposition~\ref{prop:reach-calculus} supplies the
corresponding tangent balls at second feet, with radius measured by
the source reach $d(z)$.  Exponential comparison in
\eqref{eq:diffuse-local} bounds $D$ below by a positive constant
depending only on $\delta,L_0$ within distance $10L_0$ of $X$.
Combining it with the appropriate upper bound $C_\Gamma/d(z)$
gives the asserted reach bound.  If $X$ is empty the collar assertion
has no points to consider.
\end{proof}
